% =======================================================================================
% Sommario
% =======================================================================================
\chapter*{Sommario}

L'imputazione di valori mancanti basata su \emph{Relaxed Functional Dependencies} (RFD)
affronta il problema dell'incompletezza dei dati sfruttando regole estratte dal dataset stesso.
Il costo computazionale di questi approcci cresce con il numero di regole disponibili, e una
parte di esse è ridondante: da qui l'interesse per il \emph{semantic pruning}, la selezione
delle RFD effettivamente rilevanti.

Questo lavoro studia il pruning semantico di un insieme di metadati applicato a \renuver{}, un
algoritmo di imputazione basato su RFD. Dall'analisi del bytecode si derivano tre condizioni
esatte: \Vone{}, che caratterizza le regole che possono esercitare un veto sulle imputazioni;
\Vtwo{}, che individua una classe di regole la cui potatura è dimostrabilmente priva di effetti;
\Vthree{}, che isola le regole provabilmente inerti. Da queste discende un teorema negativo: con
potatura esterna, nessuna riduzione non banale può essere contemporaneamente sicura rispetto alla
generazione dei candidati e al potere di veto.

La condizione \Vtwo{} è stata verificata empiricamente su \num{1242321} terne con zero
discrepanze, e la sua applicazione produce un output identico a quello del sistema originale su
\num{113} configurazioni su \num{113}, su quattro famiglie di dataset (\num{111} delle tre
famiglie principali e \num{2} di una famiglia indipendente). Il teorema separa ciò che è
dimostrato --- la caratterizzazione delle rimozioni certificabili --- da ciò che è misurato, e
mostra che la riduzione garantita non supera il \SI{20.9}{\percent} sui dati esaminati. Un'euristica di potatura
aggressiva, denominata \str{A5b}, riduce le regole dell'\SI{84.2}{\percent} su Bridges migliorando
al contempo il risultato di \num{69} imputazioni corrette e \num{72} errori in meno; sugli altri
dataset la stessa strategia danneggia il risultato. Il lavoro identifica il meccanismo che governa
questa differenza --- il guadagno è il prodotto della quota di imputazioni ottenute per via esatta
e del numero di imputazioni effettuate --- e ne ricava un criterio diagnostico misurabile prima di
potare. Su Glass e Restaurant una taratura più conservativa della medesima strategia non perde che
una imputazione corretta su \num{97} configurazioni. Il lavoro mostra infine che la riduzione
delle regole non implica un risparmio di tempo: la strategia che ne rimuove di più è anche quella
che rallenta di più, e il tempo dipende dalla composizione dell'insieme potato più che dalla sua
dimensione.
